\subsection{Simple and Compound Interest}


\content{
        \begin{definition} (Simple One-time Interest) The interest is only
        computed once. For example, a small loan which is payed back in full in
        three months, with some percentage of interest.
        \end{definition}
}{% presentation
\vspace{2in}
}{% nopresentation
\vspace{2in}
}{% comments
\begin{enumerate}
\item Test corrections are due today at beginning of class.
\item There will be a quiz this weekend, will open today at 1:30 PM and is due
Sunday at Midnight.
\end{enumerate}
Go over derivation of formulas. We will use $P_0$ to denote the principal (the
starting amount of a loan or account balance). 
$$ I = P_0 r $$
$$ A = P_0 + I = P_0(1+r)$$
}

\content{
        \begin{example} You lend \$300 dollars to a friend, who agrees to pay
        the amount back at the end of 3 months with 3\% interest. How much will
        you make from this? 
        \end{example}
}{% presentation
\vspace{2in}
}{% nopresentation
\vspace{2in}
}{% comments
}

\content{
        We can generalize the above idea pretty easily. 
        \begin{definition} (Simple Interest Over Time) If the interest is
        repeatedly added after a certain amount of time, denoted $t$, and the
        amount of interest does not change, this is called simple interest over
        time.
        \end{definition}
}{% presentation
\vspace{2in}
}{% nopresentation
\vspace{2in}
}{% comments
Go over the derivation of formulas $I=P_0rt$, and 
$$ A = P_0+I=P_0(1+rt).$$
}

\content{
        Most of the time, percentage rates are given as an annual percentage
        rate (APR).
        \begin{example}
        Treasury Notes (T-notes) are bonds issued by the federal government to cover its expenses.
        Suppose you obtain a \$1,000 T-note with a 4\% annual rate, paid
        semi-annually, reaching maturity in 4 years. How much interest will you earn?
        \end{example}
}{% presentation
\vspace{2in}
}{% nopresentation
\vspace{1.5in}
}{% comments
Talk about what an APR is here, since most of the examples in this chapter
will involve this.

$P_0=1000$, $r=0.02$, $t=8$ (8 6-month periods). 
}

\content{
    \begin{example} Suppose you obtain a \$1250 T-note with a 3\% APR, paid
    annually, reaching maturity in 5 years. How much interest will you earn? 
    \end{example}
}{% presentation
\vspace{2in}
}{% nopresentation
\vspace{1.5in}
}{% comments
}

\content{
        \begin{definition} (Compound Interest) When the interest is reinvested
        and can itself earn interest, this is called compound interest.
        \end{definition}
}{% presentation
\vspace{3in}
}{% nopresentation
\vspace{2in}
}{% comments
Go over derivation of formula (maybe with concrete numbers instead of
variables).
$$ P_m = P_0\left(1+\frac{r}{k}\right)^{n} $$
where $n$ is the total number of compounds, and $k$ is the number of compounds
per year. Mention that it is more common to use a formula involving the number
of years ($n=kt$). 

Mention that this is the formula we use when we deposit money into an account
and don't deposit any more, we just let it sit there. 
}

\content{
        \begin{example}
        A certificate of deposit (CD) is a savings instrument that many banks offer. It usually gives a
        higher interest rate, but you cannot access your investment for a specified length of time.
        Suppose you deposit \$3000 in a CD paying 6\% interest, compounded monthly. How much
        will you have in the account after 20 years?
        \end{example}
}{% presentation
\vspace{3in}
}{% nopresentation
\vspace{2in}
}{% comments
}

\content{
        \newpage
        \begin{example} 
        You open a checking account with an APR of 3.5\% by depositing \$1500 at
        your local bank's branch office. Assuming you don't deposit more money
        into the account, how much will be in the account in 12 years if
        interest is compounded once a year? What if interest is compounded each
        quarter? How about each month? 
        \end{example}
}{% presentation
\vspace{4in}
}{% nopresentation
\vspace{3in}
}{% comments
}
